  \documentclass[12pt]{article}

  \usepackage{sbc-template}
  \usepackage{graphicx,url}
  \usepackage[utf8]{inputenc}
  \usepackage[brazil]{babel}
  \usepackage{xcolor}

  \sloppy

  \title{Pipeline Baseada em LLM para Análise Automatizada de Editais de Fomento}

  \author{Gabriel Lopes de Albuquerque, Augusto César Ferreira de Miranda Oliveira,\\ Francisco de Assis Carlos Filho}

  \address{Universidade de Pernambuco (UPE) -- Campus Surubim\\
    Surubim -- PE -- Brasil
    \email{gabriel.lalbuquerque@upe.br, augusto.oliveira@upe.br, francisco.assis@upe.br}
  }

  \begin{document}

  \maketitle

  \begin{abstract}
    Manually monitoring public funding calls, spread across portals with different formats, is costly and makes it hard for each call to reach the researchers in its area in time. This work proposes a pipeline that collects calls from FACEPE, CNPq, CAPES, and FINEP, extracts the text of the PDFs (using OCR when needed), and uses the Gemini 2.5-Flash LLM, guided by a structured prompt, to produce data that route each call to the interested faculty. We present the problem, the hypotheses of automatic targeting (H1) and extraction reliability (H2), and the theoretical background of the proposal.
  \end{abstract}

  \begin{resumo}
    O acompanhamento manual de editais de fomento público, distribuídos em portais com formatos distintos, é custoso e dificulta que cada edital chegue a tempo aos pesquisadores de sua área. Este trabalho propõe uma pipeline que coleta editais da FACEPE, CNPq, CAPES e FINEP, extrai o texto dos PDFs (com OCR quando necessário) e usa o LLM Gemini 2.5-Flash, orientado por prompt estruturado, para gerar dados que direcionam cada edital aos docentes interessados. São apresentados o problema, as hipóteses de direcionamento automático (H1) e de confiabilidade da extração (H2) e a fundamentação teórica da proposta.
  \end{resumo}

  \noindent\textbf{Palavras-chave:} web scraping, extração de informação de documentos, large language models, prompt engineering, data pipeline, editais públicos

  \noindent\textbf{Keywords:} web scraping, document information extraction, large language models, prompt engineering, data pipeline, public funding calls

  \section{Introdução}

  Instituições de pesquisa e pesquisadores individuais dependem do acompanhamento contínuo de editais de fomento público, isto é, chamadas de financiamento, bolsas e programas de apoio, publicados por agências como a Fundação de Amparo à Ciência e Tecnologia do Estado de Pernambuco (FACEPE), o Conselho Nacional de Desenvolvimento Científico e Tecnológico (CNPq), a Coordenação de Aperfeiçoamento de Pessoal de Nível Superior (CAPES) e a Financiadora de Estudos e Projetos (FINEP). Cada uma dessas instituições publica seus editais em portais próprios, com estruturas de página, frequência de atualização e formatos de documento distintos entre si: em geral, arquivos PDF sem estrutura semântica clara \cite{b4} e, com frequência, sem disponibilização via API, situação semelhante à descrita por Krc e Schlaack [2023] para fontes governamentais que publicam seu conteúdo em páginas web sem oferecer os documentos por meio de API ou de outras tecnologias de web semântica.

  Nesse cenário, o acompanhamento manual de editais distribuídos em múltiplos portais institucionais é uma tarefa repetitiva, sujeita a atrasos e a erros de leitura. A literatura registra esse custo em contextos próximos: a extração manual de dados de documentos é lenta, custosa e sujeita a erros \cite{b4}, pode consumir várias horas de trabalho por dia \cite{b6} e, no monitoramento contínuo de fontes governamentais, chegou a exigir a dedicação integral de um profissional \cite{b3}. No próprio domínio de editais, Holanda \emph{et al.} [2026] relatam que uma universidade pública brasileira ainda reescreve manualmente seus editais, com esforço repetitivo e maior tempo de publicação. O problema é agravado pelo fato de que parte desses PDFs sequer possui camada de texto extraível (documentos escaneados ou gerados por impressão para PDF a partir de visualizadores que bloqueiam a seleção de texto), caso em que a qualidade do reconhecimento óptico de caracteres (OCR) passa a comprometer a precisão da extração \cite{b9}. O resultado é o risco concreto de perda de prazos de submissão e de decisões de elegibilidade tomadas a partir de leitura incompleta dos critérios do edital. Ao mesmo tempo, os modelos de linguagem de grande porte (LLMs) passaram a executar tarefas de extração de informação sem treinamento específico para cada tarefa \cite{b4,b7}, e métodos que integram OCR e LLMs mostram-se mais robustos ao ruído do reconhecimento de texto do que abordagens tradicionais \cite{b6}. É nesse contexto que se insere o desenvolvimento da pipeline proposta, que automatiza a coleta, extração e estruturação dessas informações.

  Diante desse problema, este trabalho investiga duas perguntas de pesquisa. A primeira (PP1) é se é possível, a partir dos dados estruturados extraídos pela pipeline (áreas de interesse, segmentos e público-alvo), direcionar automaticamente cada edital aos pesquisadores, departamentos ou grupos de pesquisa mais aderentes de sua área, de modo que esses interessados possam se organizar com antecedência e a instituição venha a captar mais recursos e ganhar mais visibilidade junto às agências de fomento. A segunda (PP2) é se é possível construir uma pipeline automatizada, baseada em um LLM, capaz de extrair e estruturar com confiabilidade suficiente os campos essenciais de um edital de fomento público, incluindo documentos sem camada de texto extraível, de modo a viabilizar, como pré-condição técnica, o direcionamento automático investigado na PP1, reduzindo ao mesmo tempo o esforço manual de acompanhamento.

  A relevância deste trabalho reside no enfrentamento da assimetria de informação e do esforço manual que historicamente limitam a captação de recursos acadêmicos. Atualmente, o acompanhamento descentralizado e intempestivo de editais resulta na perda de oportunidades e no engajamento reativo de pesquisadores. Para solucionar esse problema, propõe-se uma abordagem de caráter quantitativo articulada em torno de duas hipóteses interdependentes. A Hipótese 1 mensura o impacto final da ferramenta ao validar a eficácia do direcionamento automatizado e personalizado de oportunidades aos grupos de pesquisa e departamentos, e representa o diferencial estratégico e a contribuição principal da pesquisa: a proatividade e a assertividade na gestão de oportunidades institucionais. Por sua vez, a Hipótese 2 estabelece a base metodológica ao avaliar experimentalmente a acurácia da extração automatizada de dados de editais, pré-requisito técnico para a viabilidade do sistema. Dessa forma, a validação quantitativa da sustentação técnica (H2) garante a precisão necessária para que o direcionamento proposto na H1 seja alcançado.

  A Hipótese 1 (H1) é a de que o direcionamento automático de cada edital aos interessados certos de sua área, obtido pelo cruzamento entre os campos de classificação já extraídos pela pipeline (áreas de interesse, segmentos e público-alvo) e um cadastro de perfis institucionais (docentes, departamentos, grupos de pesquisa), permite que esses interessados se organizem com antecedência suficiente para preparar e submeter propostas mais competitivas. Como consequência, espera-se um aumento mensurável no volume de recursos de fomento captados pela instituição e ganho de visibilidade junto às agências financiadoras, em comparação com o cenário de acompanhamento manual e não direcionado.

  A Hipótese 2 (H2) é a de que uma pipeline que combina coleta específica por fonte (web scraping e integração via API), extração de texto com fallback de reconhecimento óptico de caracteres (OCR) e um LLM orientado por prompt estruturado com vocabulário controlado (dado que a qualidade das respostas de um LLM depende diretamente da estrutura e da clareza do prompt \cite{b1,b5}) é capaz de extrair, de forma consistente, os campos essenciais de um edital (público-alvo, critérios de elegibilidade, cronograma e prazo de submissão), atingindo precisão e recall por campo suficientes para sustentar o direcionamento automático da H1 e reduzindo também o esforço manual de monitoramento, sem comprometer significativamente a confiabilidade das informações extraídas. Trata-se de um ganho relevante, porém secundário diante do objetivo de direcionamento e captação de recursos investigado pela H1.

  Assim, o objetivo geral deste trabalho é desenvolver e avaliar uma pipeline automatizada capaz de coletar, extrair e estruturar, por meio de um LLM, informações de editais de fomento público publicados por múltiplas instituições, disponibilizando os resultados de forma consultável e com notificação proativa de prazos, e direcionar automaticamente cada edital aos interessados mais aderentes de sua área a partir dos campos de classificação extraídos (H1), deixando para trabalhos futuros a avaliação quase-experimental dos efeitos desse direcionamento. Para alcançá-lo, foram definidos os seguintes objetivos específicos:
  \begin{itemize}
  \item Implementar coletores (scrapers) específicos para cada fonte de dados monitorada (FACEPE, CNPq, CAPES e FINEP), respeitando as particularidades estruturais de cada portal.
  \item Extrair o conteúdo textual de documentos em PDF, incluindo casos sem camada de texto nativa, por meio de reconhecimento óptico de caracteres (OCR).
  \item Utilizar um LLM orientado por prompt estruturado para transformar o texto extraído em dados estruturados, com vocabulário controlado para os campos de classificação (áreas de interesse e segmentos).
  \item Persistir os dados extraídos em um banco de dados não relacional (MongoDB), com controle de duplicidade por documento.
  \item Notificar automaticamente, por e-mail, a existência de novos editais processados e a proximidade de prazos de submissão.
  \item Disponibilizar uma interface web para que pesquisadores marquem editais de interesse, integrada à mesma base de dados utilizada pela pipeline.
  \item Cruzar os campos de classificação extraídos de cada edital com o cadastro de preferências de docentes, direcionando automaticamente os editais novos aos interessados aderentes de sua área.
  \end{itemize}

  O restante deste documento, correspondente à primeira entrega do trabalho, está organizado da seguinte forma: a Seção~\ref{sec:fundamentacao-teorica-e-trabalhos-relacionados} apresenta os conceitos fundamentais e os trabalhos relacionados à extração de informação de documentos com LLMs e à engenharia de prompt. A metodologia, os resultados, a discussão e a conclusão serão apresentados nas entregas seguintes.

  \section{Fundamentação Teórica e Trabalhos Relacionados}\label{sec:fundamentacao-teorica-e-trabalhos-relacionados}

  \subsection{Conceitos Fundamentais}\label{sec:conceitos-fundamentais}
  Os conceitos apresentados nesta subseção acompanham as etapas do próprio pipeline do sistema proposto: a coleta dos documentos (\emph{web scraping}), a estrutura que organiza essa coleta em um fluxo reprodutível (\emph{data pipeline}), a tarefa de converter o conteúdo desses documentos em dados estruturados (\emph{Document Information Extraction}), a tecnologia empregada para realizar essa tarefa (LLMs) e a técnica usada para controlar seu comportamento (\emph{prompt engineering}).

  \textbf{Web scraping} é a técnica de extração automatizada de dados a partir de páginas web, geralmente por meio da análise de sua estrutura HTML \cite{b1}. É a abordagem adotada para três das quatro fontes do sistema proposto, cujos portais publicam os editais em páginas HTML e não disponibilizam os documentos por meio de uma API. Essa etapa de coleta é a primeira de um fluxo mais amplo, descrito a seguir como \emph{data pipeline}.

  Uma \textbf{data pipeline} (ou fluxo ETL, sigla de \emph{Extract, Transform, Load}) é uma sequência automatizada de etapas que move dados de uma fonte de origem até um destino final, aplicando transformações intermediárias \cite{b3}. No caso do sistema proposto, essas etapas correspondem à coleta (\emph{extract}), à extração de texto e estruturação via LLM (\emph{transform}) e à persistência em banco de dados (\emph{load}). A etapa de transformação, em particular, corresponde à tarefa de extração de informação de documentos, discutida a seguir.

  \textbf{Document Information Extraction} (ou \emph{PDF Data Extraction}) refere-se à tarefa de converter o conteúdo de um documento, tipicamente um PDF, em dados estruturados e semanticamente significativos. Moreira-Filho \emph{et al.} [2025] observam que os PDFs não possuem uma estrutura semântica clara, o que dificulta a identificação de elementos como títulos, tabelas e seções, e tratam a extração estruturada como um problema de reconhecimento de entidades nomeadas (NER) no qual os LLMs se mostram eficientes graças à sua capacidade de compreensão contextual. É justamente essa tecnologia, os LLMs, que se descreve na sequência.

  \textbf{Large Language Models (LLMs)} são modelos de linguagem baseados em técnicas de processamento de linguagem natural, treinados por aprendizado auto-supervisionado em um grande corpus de texto \cite{b7}. Sua escala cresceu de forma acentuada nos últimos anos, de centenas de milhões a mais de um trilhão de parâmetros \cite{b7}, e esses modelos apresentam forte capacidade de generalização, executando tarefas variadas, inclusive a extração de informação a partir de texto, sem treinamento específico para cada tarefa \cite{b7}. Em vez de atualizar os parâmetros do modelo, é possível adaptá-lo a uma tarefa nova apenas por meio do prompt \cite{b5}, o que permite que um LLM siga instruções detalhadas em linguagem natural para estruturar um edital a partir apenas do texto e do prompt fornecidos. Controlar esse comportamento por meio do prompt é, por sua vez, o objeto da técnica descrita a seguir.

  \textbf{Prompt engineering} é definida como o processo de criar e ajustar instruções específicas (o \emph{prompt}) para guiar o comportamento de um LLM e obter respostas mais precisas e consistentes \cite{b5,b1}. Entre as estratégias mais estabelecidas estão o \emph{zero-shot prompting} (nenhum exemplo fornecido), o \emph{few-shot prompting} (poucos exemplos fornecidos no próprio prompt) e o \emph{chain-of-thought (CoT) prompting} (instrução para que o modelo explicite passos intermediários de raciocínio antes da resposta final). O prompt utilizado pelo analisador do sistema proposto combina uma instrução de papel (``atue como um analista de editais''), restrições explícitas de formato de saída (JSON) e um vocabulário controlado para dois dos campos extraídos, numa abordagem próxima da técnica de \emph{role prompting} combinada a \emph{few-shot prompting} descrita por Blázquez-Ochando, Prieto-Gutiérrez e Ovalle-Perandones [2025].

  \subsection{Trabalhos Relacionados}

  \subsubsection{Método de Seleção dos Trabalhos}
  Os trabalhos discutidos nesta seção foram reunidos por um mapeamento sistemático da literatura recente em torno dos três eixos temáticos mais próximos do sistema proposto: extração de informação de documentos com LLMs, engenharia de prompt aplicada a tarefas estruturadas e pipelines automatizadas de coleta de documentos. A busca foi conduzida no Google Scholar, pela sua ampla cobertura, e complementada por consulta direta a repositórios e bibliotecas digitais (\emph{arXiv}, IEEE Xplore, Wiley Online Library e periódicos da MDPI) e aos anais do Simpósio Brasileiro de Sistemas de Informação (SBSI), combinando os termos ``web scraping'', ``data pipeline'', ``ETL'', ``document information extraction'', ``PDF data extraction'', ``document understanding'', ``large language models'', ``semantic extraction'', ``document classification'' e ``prompt engineering''. Essa varredura inicial identificou cerca de 20 trabalhos aparentemente aderentes ao escopo. Uma primeira etapa de triagem, por leitura de título, resumo e palavras-chave, manteve os trabalhos publicados entre 2023 e 2025 com aplicação prática ou validação experimental sobre uso de LLMs, extração de dados, automação de documentos ou \emph{prompt engineering}, e excluiu duplicatas, trabalhos de abordagem puramente teórica e os voltados à extração de fontes já estruturadas (como APIs), reduzindo o conjunto a cerca de 10 trabalhos. A segunda etapa, de leitura integral, selecionou 6 trabalhos \cite{b1,b2,b3,b4,b5,b6} como base principal da fundamentação. A esse núcleo somaram-se, posteriormente, um \emph{survey} \cite{b7} usado apenas como referência de enquadramento teórico geral sobre características de LLMs (Seção~\ref{sec:conceitos-fundamentais}) e os dois trabalhos do SBSI \cite{b8,b9}, localizados por busca manual nos anais do evento, totalizando os nove trabalhos aqui discutidos. Sahoo \emph{et al.} [2024] foi localizado pelo próprio web scraper do projeto; os demais, por busca manual.

  \subsubsection{Extração de Dados de Literatura Científica com LLMs e KNIME}
  Moreira-Filho \emph{et al.} [2025] partem da constatação de que o crescimento da literatura científica (mais de 2,5 milhões de artigos publicados por ano) torna a extração manual de dados inviável em escala, e descrevem um \emph{workflow}, construído na plataforma \emph{low-code} KNIME, que integra ferramentas de \emph{parsing} de documentos com LLMs para extrair variáveis de publicações científicas e de PDFs em geral, tratando a tarefa como um problema de reconhecimento de entidades nomeadas (NER) no qual o LLM se mostra mais eficiente que abordagens tradicionais graças à sua capacidade de compreensão contextual. O sistema opera em dois modos: um modo texto, que extrai diretamente texto e tabelas, e um modo imagem, que envia páginas renderizadas a um LLM multimodal. A engenharia de \emph{prompt} é apontada pelos autores como fundamental para garantir precisão e consistência e evitar alucinações, com a temperatura fixada em zero para resultados determinísticos e a saída sempre em JSON. Os autores reportam 81,14\% de acurácia no modo texto para publicações científicas e até 98,54\% no modo imagem para PDFs gerais, com Claude 3.5 Sonnet e GPT-4o obtendo o melhor desempenho e a menor incidência de alucinação entre os modelos testados. É o trabalho mais próximo estruturalmente do sistema proposto entre os consultados, e as decisões de \emph{prompt} descritas (restrição de formato de saída e temperatura zero) são próximas das adotadas no analisador do sistema proposto.


  \subsubsection{Pipeline de Metadados em Escala a partir de Web Scraping}
  Krc e Schlaack [2023] relatam a construção de uma pipeline para a biblioteca digital FRASER (Federal Reserve Bank of St. Louis), que automatiza a coleta, a transformação e a ingestão de documentos e metadados de fontes governamentais dos Estados Unidos que, segundo os autores, ``apresentam seu conteúdo em páginas web, mas não disponibilizam os metadados e links de documentos via API ou outras tecnologias de web semântica'' \cite{b3}. O processo, originalmente totalmente manual e não projetado para automação, cresceu a ponto de o monitoramento de alertas e a coleta, descrição e ingestão manuais se tornarem, já em 2021, uma posição de dedicação integral. A solução resultante é um programa em Python executado diariamente que faz o \emph{scraping} de uma página, verifica se os itens capturados já existem no repositório e, em caso negativo, transforma e ingere os metadados e o arquivo e envia uma notificação por e-mail; a arquitetura separa o processamento de dados do gerenciamento de dados, define o que coletar por meio de arquivos JSON (\emph{sitemaps} com seletores) e aciona uma função \emph{serverless} (AWS Lambda) a partir de um \emph{endpoint} HTTP. A deduplicação é feita pela comparação do nome de arquivo construído contra os arquivos já armazenados em um \emph{bucket} AWS S3, e uma etapa manual de conferência e aprovação precede a ingestão final. Os autores reportam que o tempo ativo para adicionar novos registros foi reduzido a poucos cliques e que os alertas de falso positivo foram reduzidos a zero. É o trabalho mais próximo do sistema proposto em termos de escopo (coleta contínua e em escala de documentos governamentais heterogêneos), mas não emprega um LLM em nenhuma etapa: a curadoria final fica a cargo de uma conferência humana.

  \subsubsection{Simplificação Automática de Editais com LLM}
  Holanda \emph{et al.} [2026], em trabalho publicado no SBSI, avaliam uma pipeline baseada em LLM que gera, em modo \emph{zero-shot}, versões em linguagem simples de editais e comunicados administrativos de uma universidade pública, hoje reescritos manualmente pela instituição, com custo de tempo e retrabalho. A etapa de entrada coincide com a do sistema proposto: os editais publicados on-line são extraídos e convertidos de PDF para texto antes de irem para o \emph{prompt}. Os autores usam um \emph{prompt} fixo, fragmentam o texto em blocos de até 4.000 caracteres e fixam a temperatura em zero para garantir replicabilidade. A avaliação combina índices estatísticos de legibilidade e métricas morfossintáticas do projeto NILC-Metrix com um questionário aplicado por conveniência a docentes e estudantes que leem ou submetem editais ao menos uma vez por ano. Dois resultados são especialmente relevantes para este trabalho: na comparação entre modelos, o Gemini Flash superou o Gemini Pro em todas as métricas de legibilidade, e entre os modelos abertos o Qwen2.5:14b foi o que mais se aproximou; e vários modelos alucinaram durante a geração (produzindo textos ininteligíveis, repetindo o \emph{prompt} ou respondendo em outro idioma), o que exigiu revisão manual de todas as saídas e o registro do número de documentos alucinados por modelo. O estudo reforça, para o contexto brasileiro de documentos administrativos públicos, tanto a escolha de um modelo da linha ``Flash'' quanto a necessidade de uma camada de validação sobre a saída do LLM, ambas presentes no sistema proposto.

  \subsection{Comparação dos Trabalhos Relacionados}
  Nenhum dos trabalhos consultados trata especificamente de editais de fomento público em português, nem combina, em uma única pipeline de produção, os quatro elementos centrais do sistema proposto: (i) coleta contínua e automatizada de múltiplas fontes governamentais heterogêneas; (ii) extração de texto com \emph{fallback} de OCR para documentos sem camada de texto nativa; (iii) estruturação via LLM com vocabulário controlado; e (iv) notificação proativa sensível a prazos. Krc e Schlaack [2023] compartilham o desafio de escala e heterogeneidade de fontes governamentais, mas não empregam um LLM em nenhuma etapa, apoiando-se em conferência humana. Moreira-Filho \emph{et al.} [2025] e Chen \emph{et al.} [2025] empregam LLMs para extração estruturada e reportam métricas quantitativas de acurácia sobre conjuntos de teste, rigor de avaliação que o sistema proposto ainda não possui de forma formal; isso configura uma lacuna concreta e um direcionamento para trabalhos futuros. Blázquez-Ochando, Prieto-Gutiérrez e Ovalle-Perandones [2025] usam o LLM para gerar o código de coleta, e não para interpretar o conteúdo coletado, e Fahrudin \emph{et al.} [2025] usam o LLM apenas de forma indireta, por meio de \emph{embeddings} de similaridade para ranqueamento de relevância. Os dois trabalhos do SBSI aproximam-se pelo domínio: Holanda \emph{et al.} [2026] aplicam um LLM a editais de uma universidade pública brasileira, ainda que para simplificação e não para extração estruturada, e Utyiama \emph{et al.} [2026] adotam uma arquitetura de extração e validação de campos de documentos em PDF próxima da do sistema proposto, porém no domínio fiscal e a partir de documentos já recebidos pelo fluxo de trabalho da organização, sem coleta automatizada das fontes.

  

  \begin{thebibliography}{00}
  \bibitem[Blázquez-Ochando et al. 2025]{b1} Blázquez-Ochando, M., Prieto-Gutiérrez, J. J., and Ovalle-Perandones, M. A. (2025). Prompt engineering for bibliographic web-scraping. \emph{Scientometrics}, 130:3433--3453.
  \bibitem[Chen et al. 2025]{b6} Chen, L.-C. et al. (2025). Evaluation of prompt engineering on the performance of a large language model in document information extraction. \emph{Electronics}, 14(11):2145.
  \bibitem[Fahrudin et al. 2025]{b2} Fahrudin, T. M., Funabiki, N., Brata, K. C., Naing, I., Aung, S. T., Muhaimin, A., and Prasetya, D. A. (2025). An improved reference paper collection system using web scraping with three enhancements. \emph{Future Internet}, 17(5):195.
  \bibitem[Holanda et al. 2026]{b8} Holanda, J. P., Magalhães, R., Almendra, C., and do Rêgo, L. G. C. (2026). Simplifying administrative texts for plain language using LLM: a model comparative analysis. In \emph{Anais do XXII Simpósio Brasileiro de Sistemas de Informação (SBSI)}, pages 931--947, Porto Alegre. SBC.
  \bibitem[Krc and Schlaack 2023]{b3} Krc, M. and Schlaack, A. O. (2023). Pipeline or pipe dream: Building a scaled automated metadata creation and ingest workflow using web scraping tools. \emph{Code4Lib Journal}, (58).
  \bibitem[Moreira-Filho et al. 2025]{b4} Moreira-Filho, J. T. et al. (2025). Automating data extraction from scientific literature and general PDF files using large language models and KNIME: An application in toxicology. \emph{WIREs Computational Molecular Science}, 15:e70047.
  \bibitem[Sahoo et al. 2024]{b5} Sahoo, P., Singh, A. K., Saha, S., Jain, V., Mondal, S., and Chadha, A. (2024). A systematic survey of prompt engineering in large language models: Techniques and applications. \emph{arXiv preprint arXiv:2402.07927}.
  \bibitem[Utyiama et al. 2026]{b9} Utyiama, D. M. da S., Moraes, O. S. C., Guimarães, F. A. S., Ramos, P. B. S., and Marques, L. C. (2026). Smart OCR: a vertical AI agent for financial document validation, an exploratory case study. In \emph{Anais do XXII Simpósio Brasileiro de Sistemas de Informação (SBSI)}, pages 948--967, Porto Alegre. SBC.
  \bibitem[Xu et al. 2024]{b7} Xu, H., Gan, W., Qi, Z., Wu, J., and Yu, P. S. (2024). Large language models for education: A survey. \emph{arXiv preprint arXiv:2405.13001}.
  \end{thebibliography}

  \end{document}
